% ============================================================
% 2. BACKGROUND / RELATED WORK
% Erfan owns 2.1 (JEPA/VL-JEPA) -- sized stub.
% Parsa: 2.2 (gaze-conditioning) + 2.3 (project origin, from EARLY_WORK_SUMMARY.md).
% 2.3 REWRITTEN 2026-08-13: real origin = 5-phase plan -> confounded 2.81/4.27
% pilot. Context-selection thread demoted to one sentence.
% ============================================================
\section{Background and related work}

% ------------------------------------------------------------
\subsection{Joint-embedding predictive architectures for video}
\label{subsec:bg-jepa}
% -----------------------------------------------------------------------------
% ERFAN: the JEPA lineage that sets up our method. Suggested coverage:
%   - I-JEPA (image, latent prediction vs pixel reconstruction) \citep{assran2023ijepa}
%   - V-JEPA / V-JEPA 2 / 2.1 (video feature prediction; frozen encoder+predictor)
%     \citep{bardes2024vjepa, assran2025vjepa2, murlabadia2026vjepa21}
%   - DINO-WM (world models on frozen features) \citep{zhou2024dinowm}
%   - VL-JEPA / language-aligned latent prediction -- YOUR section, tie to horizon.
%   - EgoAgent (joint predictive agent model, egocentric) \citep{chen2025egoagent} -- adjacent framing.
%   - the AC (action-conditioned) predictor as the thing we re-purpose for HUMAN signals.
% Keep it tight; \S\ref{subsec:method-setup} carries the setup detail.
% -----------------------------------------------------------------------------
\paragraph{Latent prediction.}
The lineage this work builds on replaces pixel reconstruction with prediction in
representation space. I-JEPA~\citep{assran2023ijepa} established the pattern for
images: from a context block, predict the \emph{embeddings} of masked target
blocks rather than their pixels, so the objective is never forced to model detail
that carries no semantic content. V-JEPA~\citep{bardes2024vjepa} carried the idea
to video, and V-JEPA~2~\citep{assran2025vjepa2} scaled it. Its three components
fix the vocabulary we use throughout: a student encoder that sees the visible
patches, a teacher encoder updated by exponential moving average that supplies
target embeddings under a stop-gradient, and a \emph{predictor} that maps the
student's context embeddings plus positional queries for the masked positions
onto what the teacher would emit there. The entire training signal is a distance
between predicted and target latents; nothing is ever rendered to pixels.
V-JEPA~2.1~\citep{murlabadia2026vjepa21} revises the recipe toward denser
features, and is the backbone of both frozen-probe systems that currently lead
EK100 anticipation (\S\ref{subsec:bg-gaze}).

\paragraph{Frozen features as a world model.}
The encoder/predictor split matters because the two can be used separately.
DINO-WM~\citep{zhou2024dinowm} builds a world model on top of frozen visual
features and plans by optimising action sequences against predicted future
latents, which is evidence that a frozen representation plus a learned latent
predictor is already enough to support prediction-based control. Our protocol
rests on the same assumption in a weaker form: we never fine-tune the
representation, and read anticipation out of it with a deliberately small probe
(\S\ref{subsec:method-setup}).

\paragraph{The action-conditioned predictor, and what we re-purpose.}
V-JEPA~2 also ships an action-conditioned variant (V-JEPA~2-AC), in which the
predictor is additionally conditioned on robot end-effector action and state,
each encoded to a token and concatenated as a peer alongside the visual tokens
\emph{inside} the predictor. That token layout is precisely the interface this
paper re-purposes: we substitute human gaze and hand pose for the robot's action
and state and change nothing else about the pathway
(\S\ref{subsec:method-weak}). EgoAgent~\citep{chen2025egoagent} makes an adjacent
move in the egocentric setting, learning representation, prediction and action
jointly rather than treating them as separate stages.

\paragraph{Language-aligned latent prediction.}
A final branch of the lineage moves the \emph{target} of latent prediction from
appearance toward meaning. VL-JEPA~\citep{vljepa2025} predicts the embedding of
an answer rather than generating it token by token, so that two paraphrases of
the same content land near each other instead of being orthogonal in token space,
and the semantics are available without waiting for a decode. Its relevance here
is directional rather than architectural: it demonstrates that the JEPA objective
survives a change of target, which is exactly the move our own results argue for
in place of input-level conditioning --- supervising the representation toward
goal- or intention-level targets rather than concatenating a signal at the input
(\S\ref{sec:conclusion}).

% ------------------------------------------------------------
\subsection{Conditioning prediction on behavioral signals}
\label{subsec:bg-gaze}

A parallel line of work asks how gaze and body signals should enter a predictive
model, and it divides usefully by where the signal is injected. At one end sit
input concatenation methods, the peer-token approach this paper diagnoses. At the
other, signals supervise the representation during training and are discarded at
inference. GABRIL \citep{banayeeanzade2025gabril} regularises an
imitation-learning representation toward human gaze to mitigate causal confusion;
the gaze-regularised vision-language models of \citet{pani2026gazereg} align model
attention to gaze via a regularisation term and report improved future-event
prediction without gaze at inference; and \citet{royfernando2022abstractgoal}
model an abstract goal distribution and a goal-consistency objective for
next-action selection. A related line conditions anticipation explicitly on
intention or goal: the intention-conditioned VAE of \citet{mascaro2023icvae}, the
goal-conditioning isolated by AntGPT \citep{zhao2024antgpt}, the intention-guided
cognitive reasoning of \citet{chu2026intention}, and the gaze-guided,
intention-conditioned graph network of \citet{ozdel2024gazegnn} all report that an
intention or goal signal improves long-term anticipation, a pattern our own
results echo in pointing away from input-level conditioning and toward goal-level
targets. Anticipating the behavioral signal itself has also been studied
directly, as in the audio-visual egocentric gaze anticipation of
\citet{lai2024listentolook}. Between these, GazeQwen \citep{pham2026gazeqwen}
modulates a frozen video-language model with gaze through a lightweight resampler
that adds residuals at decoder layers, and argues that the obstacle is a missing
learned pathway from gaze into the model's internal representations rather than a
deficiency of the signal. That framing organises our own investigation: we adapt
GazeQwen's mechanisms as the template for a stronger pathway, but test whether
such a pathway could help before building it (\S\ref{subsec:method-strong}).

On the benchmark itself, the two strongest published EK100 anticipation systems
are both frozen-V-JEPA-2.1 attentive probes \citep{chu2026jfaa, wang2026tapjepa}.
Neither conditions on behavioral signals, yet both split their probes into
separate verb and noun heads on the reasoning that verbs track motion and nouns
track object appearance, a split that reappears in our own results
(\S\ref{subsec:dissociation}). They mark the vision-only ceiling of the
frozen-probe regime, and they leave our question, whether behavioral conditioning
adds anything on top, untested.

% ------------------------------------------------------------
\subsection{Project origin}
\label{subsec:bg-origin}

This study grew out of an earlier attempt to carry V-JEPA~2's action-conditioned
predictor from robot manipulation to egocentric human video, replacing the
robot's end-effector action and state tokens with human gaze and hand-pose
tokens. (An even earlier exploratory step had used gaze only to bias
\emph{context selection} for masked reconstruction, and found no benefit at high
coverage; it did not lead anywhere and we mention it only for completeness.) The
conditioning effort was organised as a five-stage plan: reproduce the published
EK100 anticipation baseline, train an own probe to own that baseline, map the
action-conditioned predictor's architecture, post-train the predictor on HD-EPIC
with gaze and hand tokens, and re-probe on EK100 to measure the change.

The plan did not survive contact with a deadline, and the first behavioral
comparison was confounded by design. The from-scratch behavioral predictor
trained on HD-EPIC was a ViT-L model, but the only fine-tuned
behavioral checkpoint available in time was the ViT-G model from the project's
own conditioning track, trained separately and with a differently-shaped
(three-dimensional) gaze input. Rather than wait, the comparison
was run across encoders, knowingly unfairly: a ViT-L vision baseline reached
$4.27\%$ action R@5 on EK100 while the ViT-G behavioral arm reached $2.81\%$, with
the behavioral arm also behind on verb and noun recall, feature silhouette, and
nearest-neighbour retrieval. Both sides were trained for a single probe epoch and
were underfit. The number is therefore directional only, and we treat it as a
confounded pilot rather than a result: the encoders differ, the predictors come
from different training runs, and gaze is masked at EK100 evaluation. What the
pilot did establish was the question this paper takes up. Behavioral conditioning
had not helped, and it was not yet clear whether the fault lay with the injection
pathway, the confounds, or the signal itself. The controlled study that follows
was built to separate those explanations.
% NOTE: origin sourced from EARLY_WORK_SUMMARY.md (reconstructed from ocean-cluster
% session transcripts + live CLAUDE.md/eval_report.md). 4.27 = Phase 2 ViT-L
% baseline; 2.81 = Phase 5 ViT-G behavioral (colleague checkpoint). Deliberately
% confounded, directional only -- do NOT upgrade to a headline result.